% © 2026 Ing. David Jaroš — CC BY-NC-ND 4.0
%
% This work is licensed under a Creative Commons Attribution-NonCommercial-NoDerivatives
% 4.0 International License (CC BY-NC-ND 4.0).
%
% License History: Earlier drafts (up to v0.3) were released under CC BY 4.0.
% From v0.4 onward, all material is released under CC BY-NC-ND 4.0 to protect
% the integrity of the theoretical work during ongoing academic development.
%
% See LICENSE.md for full license text.
%
% Track: T3_ALPHA — Fine Structure Constant
% Purpose: Gap G137-B — Mellin insertion B (Priority 1)
% Predecessor: chowla_selberg_B_derivation.tex (Steps 1–7)
% Status: Gap G137-B NARROWED — Mellin insertion step pending [OPEN/MC]
% Date: 2026-05-14

\documentclass[12pt,a4paper]{article}
\usepackage[a4paper,margin=2.5cm]{geometry}
\usepackage{amsmath,amssymb,amsthm,booktabs,tcolorbox}
\tcbuselibrary{skins,breakable}
\newtcolorbox{statusbox}[1][]{%
  colback=gray!8!white,colframe=black!60,arc=3pt,boxrule=1pt,
  title=Status Summary,#1}
\newtheorem{theorem}{Theorem}[section]
\newtheorem{lemma}[theorem]{Lemma}
\newtheorem{proposition}[theorem]{Proposition}
\newtheorem{corollary}[theorem]{Corollary}

\title{Gap G137-B: Mellin Insertion for Coefficient $B$\\
       on $\mathbb{Z}^2$ at the Self-Dual Point}
\author{Ing.\ David Jaroš}
\date{2026-05-14}

\begin{document}
\maketitle

\begin{abstract}
This note addresses the open step in Gap G137-B: deriving the coefficient
$B \approx 46.28$ in $V_{\mathrm{eff}}(n) = n^2 - Bn\ln n$ from the UBT
action via Epstein zeta functions on $T^2$, Chowla--Selberg asymptotics,
and the Mellin/spectral insertion linking the determinant to the $n\ln n$
coefficient. Steps 1--3 are carried out to level [L1 asym.] or [STD].
The Mellin insertion (Step 4) is identified as the precise remaining gap
and remains [OPEN/MC]. Gap G137-B status: \textbf{NARROWED}.
Alpha: \textbf{NOT DERIVED}.
\end{abstract}

\tableofcontents

%──────────────────────────────────────────────────────────────────────────────
\section{Setup: Epstein Zeta on $T^2$ at Self-Dual Point}
\label{sec:setup}
%──────────────────────────────────────────────────────────────────────────────

Let $T^2 = S^1_\psi \times S^1_\chi$ with equal radii $R$.
The scalar propagator operator on $T^2$ with mass shift $d^2/R^2$ is
\[
P_b = -\partial^2_\psi - \partial^2_\chi + \frac{d^2}{R^2}.
\]
Its spectral zeta function is the Epstein zeta function of $\mathbb{Z}^2$:
\[
Z_{T^2}(s;\, d^2/R^2)
= R^{-2s}\, Z_{\mathbb{Z}^2}(s;\, d^2 R^2 / \Lambda^2),
\]
where $\Lambda$ is the UV regulator and $Z_{\mathbb{Z}^2}(s; a)$
is the Epstein zeta function with parameter $a > 0$.

\subsection*{Standard values [STD]}

\begin{enumerate}
\item $Z_{\mathbb{Z}^2}(0; a) = -1$ for all $a > 0$
      (standard zeta-regularization of the constant mode sum) [STD].
\item $Z_{\mathbb{Z}^2}(1; 0) = \pi$ (the bare Epstein sum, no mass shift) [STD].
\item Self-dual point: $R\Lambda = 1$, so $Z_{T^2}(1; 0)|_{R\Lambda=1} = \pi\Lambda^2$.
\end{enumerate}

%──────────────────────────────────────────────────────────────────────────────
\section{Step 1: Zeta Derivative $Z'_{T^2}(0; d^2)$ via Chowla--Selberg}
\label{sec:step1}
%──────────────────────────────────────────────────────────────────────────────

Taking the $s$-derivative of $Z_{T^2}(s; d^2/R^2)$ at $s = 0$:
\begin{equation}
Z'_{T^2}(0;\, d^2/R^2)
= 2\ln(R)\cdot Z_{\mathbb{Z}^2}(0;\, d^2)
  + Z'_{\mathbb{Z}^2}(0;\, d^2)
= -2\ln(R) + Z'_{\mathbb{Z}^2}(0;\, d^2),
\label{eq:zt2_deriv}
\end{equation}
using $Z_{\mathbb{Z}^2}(0; a) = -1$ [STD].

\medskip\noindent
\textbf{Chowla--Selberg formula [STD]}.
For the Epstein zeta function $Z_{\mathbb{Z}^2}(s; d^2)$ of the
rank-2 lattice with parameter $d^2$, the derivative at $s = 0$ is:
\begin{equation}
Z'_{\mathbb{Z}^2}(0;\, d^2) = -4\ln|\Gamma({\tfrac{1}{4}})| + 2\ln(d) + C_0,
\label{eq:cs_deriv}
\end{equation}
where $C_0$ is an analytic constant independent of $d$, and
$\Gamma(1/4)$ arises from the functional determinant of the Laplacian on $T^2$
via the Kronecker limit formula.

\medskip\noindent
\textbf{Modular identities [STD]}:
\begin{align}
\Gamma({\tfrac{1}{4}}) &= 2\pi^{3/4}\,\eta(i),
  \label{eq:gamma14_eta}\\
\vartheta_3(0|i) &= \sqrt{2}\,\eta(i),
  \label{eq:theta3_eta}
\end{align}
where $\eta(i) \approx 0.7682$ is the Dedekind eta function at $\tau = i$,
and $\vartheta_3(0|i) \approx 1.0864$ is the Jacobi theta function [STD, NUM].

Combining \eqref{eq:zt2_deriv}--\eqref{eq:cs_deriv}:
\begin{equation}
Z'_{T^2}(0;\, d^2/R^2)
= -2\ln R - 4\ln|\Gamma({\tfrac{1}{4}})| + 2\ln(d/R) + C_0.
\label{eq:zt2_combined}
\end{equation}

\begin{statusbox}
Step 1: $Z'_{T^2}(0; d^2) = -2\ln R + Z'_{\mathbb{Z}^2}(0; d^2)$: \textbf{[STD]}\\
Chowla--Selberg $Z'_{\mathbb{Z}^2}(0; d^2) = -4\ln|\Gamma(\tfrac14)| + 2\ln d + C_0$: \textbf{[STD]}\\
$\Gamma(\tfrac14) = 2\pi^{3/4}\eta(i)$, $\vartheta_3(0|i) = \sqrt{2}\eta(i)$: \textbf{[STD]}
\end{statusbox}

%──────────────────────────────────────────────────────────────────────────────
\section{Step 2: Divisor Sum and $n\ln n$ Coefficient}
\label{sec:step2}
%──────────────────────────────────────────────────────────────────────────────

The one-loop effective action in UBT contributes to $V_{\mathrm{eff}}(n)$ via
\begin{equation}
V_{\mathrm{eff}}(n)
= -\sum_{d\mid n} Z'_{T^2}(0;\, d^2/R^2).
\label{eq:veff_divisor}
\end{equation}
Substituting \eqref{eq:zt2_combined}:
\begin{align}
-\sum_{d\mid n} Z'_{T^2}(0;\, d^2/R^2)
&= \sum_{d\mid n}\bigl[2\ln R + 4\ln|\Gamma({\tfrac14})| - 2\ln(d/R) - C_0\bigr]
\nonumber\\
&= 2\sigma_0(n)\ln R
  + 4\sigma_0(n)\ln|\Gamma({\tfrac14})|
  - 2\sum_{d\mid n}\ln d
  + 2\sigma_0(n)\ln R
  - \sigma_0(n) C_0,
\label{eq:veff_expanded}
\end{align}
where $\sigma_0(n) = \sum_{d\mid n} 1$ is the divisor counting function.

\subsection*{Asymptotic analysis}

For the coefficient of $n\ln n$ in $V_{\mathrm{eff}}$, consider the average:
\[
B = \lim_{N\to\infty}
    \frac{1}{N\ln N}
    \sum_{n=1}^N
    \biggl(-\sum_{d\mid n} Z'_{T^2}(0;\,d^2/R^2)\biggr).
\]
The key identity from Dirichlet's divisor problem [STD]:
\[
\sum_{n=1}^N \sigma_0(n) \sim N\ln N,
\qquad
\sum_{n=1}^N \sum_{d\mid n} \ln d = \sum_{d=1}^N \ln d \cdot \lfloor N/d\rfloor
\sim N\ln N.
\]
Thus the $n\ln n$ coefficient in the averaged $V_{\mathrm{eff}}$ is [L1 asym.]:
\begin{equation}
B_{\mathrm{single}} \sim 4\ln|\Gamma({\tfrac14})| + 2\ln R + \text{const}.
\label{eq:b_single}
\end{equation}

\begin{lemma}[$n\ln n$ coefficient fixed by Dirichlet asymptotic, [L1 asym.]]
\label{lem:nlogn}
The coefficient of $n\ln n$ in the one-channel averaged $V_{\mathrm{eff}}$
equals
\[
B_{\mathrm{single}} = 4\ln|\Gamma({\tfrac14})| + c_R
\]
for a constant $c_R$ depending on the $R$-normalization.
The $n\ln n$ source is uniquely the divisor asymptotic $\sum_{n\leq N}\sigma_0(n)\sim N\ln N$.
\end{lemma}

\begin{proof}
From \eqref{eq:veff_expanded}: the dominant term proportional to $\sigma_0(n)$
gives the $n\ln n$ coefficient via Dirichlet's theorem
$\sum_{n\leq N}\sigma_0(n) \sim N\ln N$.
The term $-2\sum_{d\mid n}\ln d$ contributes at the same order via
$\sum_{n\leq N}\sum_{d\mid n}\ln d \sim N\ln N$.
\end{proof}

\begin{statusbox}
Step 2: $n\ln n$ from divisor asymptotic $\sum_{n\leq N}\sigma_0(n)\sim N\ln N$: \textbf{[L1 asym.]}\\
Single-channel $B_{\mathrm{single}} = 4\ln|\Gamma(\tfrac14)| + c_R$ confirmed: \textbf{[L1 asym.]}\\
Large observed $B$-value requires $N_{\mathrm{eff}}^{3/2}$ lifting: \textbf{[NEEDS STEP 4]}
\end{statusbox}

%──────────────────────────────────────────────────────────────────────────────
\section{Step 3: $N_{\mathrm{eff}}$ Lifting and Candidate $B$}
\label{sec:step3}
%──────────────────────────────────────────────────────────────────────────────

The full UBT $V_{\mathrm{eff}}$ involves $N_{\mathrm{eff}} = 12$ complex-time
phases [L1, \texttt{su2\_twist\_neff12.tex}].  The volumetric lifting factor
$N_{\mathrm{eff}}^{3/2}$ arises from the $T^3_\psi\times S^1_\chi$ geometry
under the self-dual identification; specifically, $N_{\mathrm{eff}}^{1/2}$ from
the $T^3$ volume prefactor at self-duality and $N_{\mathrm{eff}}$ from summing
over all modes.

\subsection*{T\(^3\) factorization in terms of phase and per-phase multiplicity}
Write
\[
N_{\mathrm{eff}} = N_{\mathrm{phases}}\cdot N_{\mathrm{eff,pp}},
\qquad
N_{\mathrm{phases}}=3,\quad N_{\mathrm{eff,pp}}=4,
\]
where \(N_{\mathrm{phases}}\) comes from \(\dim_{\mathbb R}\mathrm{Im}(\mathbb H)=3\)
[\(L0\)] and \(N_{\mathrm{eff,pp}}=4\) from the SU(2) twist per phase [\(L1\)].
Then
\[
N_{\mathrm{eff}}^{3/2}
=\left(N_{\mathrm{phases}}N_{\mathrm{eff,pp}}\right)^{3/2}
=N_{\mathrm{phases}}^{3/2}N_{\mathrm{eff,pp}}^{3/2}
=3^{3/2}\cdot4^{3/2}=12^{3/2}.
\]
In this decomposition, the candidate coefficient is
\[
B
=N_{\mathrm{phases}}^{3/2}N_{\mathrm{eff,pp}}^{3/2}\,(2\eta(i))^{1/4}
=12^{3/2}(2\eta(i))^{1/4}.
\]
So the volumetric \(3/2\)-exponent is carried by the full effective multiplicity
\(N_{\mathrm{eff}}=12\), while the unresolved part remains the Mellin/heat-kernel
origin of \((2\eta(i))^{1/4}\) from \(S[\Theta]\).

The candidate formula is:
\begin{equation}
B = N_{\mathrm{eff}}^{3/2}\cdot Z_{1\mathrm{real}}(\tau=i)^{1/4},
\label{eq:B_candidate}
\end{equation}
where $Z_{1\mathrm{real}}(\tau=i) = 2\eta(i)$ is the normalization factor for
one real scalar at the self-dual point [OPEN/MC].

\subsection*{Numerical verification [NUM]}

Using the standard values $\eta(i) \approx 0.76823$,
$\vartheta_3(0|i) = \sqrt{2}\,\eta(i) \approx 1.08643$:
\begin{align}
2\eta(i) &\approx 1.5365, \label{eq:2eta}\\
(2\eta(i))^{1/4} &\approx 1.1133, \label{eq:2eta_14}\\
12^{3/2} &= 41.569, \label{eq:neff32}\\
B &= 12^{3/2}\cdot(2\eta(i))^{1/4} \approx 46.28. \label{eq:B_num}
\end{align}

This is consistent with the phenomenological value $B_{\mathrm{phenom}}\approx 46.28$
required for the prime attractor mechanism (route A\_PRIME,
\texttt{canonical/alpha/alpha\_gap\_closure\_matrix.tex}).

\subsection*{Consistency check: $N_{\mathrm{eff}} = 12$ versus $N_{\mathrm{phases}} = 3$}

Using $N_{\mathrm{phases}} = 3$ instead gives
$3^{3/2}\cdot(2\eta(i))^{1/4} \approx 5.196\times 1.113 \approx 5.79$,
far below $B_{\mathrm{phenom}}$.  Hence the SU(2)-twist $N_{\mathrm{eff}} = 12$
is essential [L1, \texttt{su2\_twist\_neff12.tex}].

\begin{statusbox}
Step 3: $N_{\mathrm{eff}} = 12$ [L1], $N_{\mathrm{eff}}^{3/2} = 41.57$: \textbf{[L1]}\\
$N_{\mathrm{eff}}^{3/2}=N_{\mathrm{phases}}^{3/2}N_{\mathrm{eff,pp}}^{3/2}$ with $N_{\mathrm{phases}}=3$, $N_{\mathrm{eff,pp}}=4$: \textbf{[L1/L0 decomposition]}\\
$Z_{1\mathrm{real}} = 2\eta(i)$: \textbf{[OPEN/MC]} (set by hand; not derived from NS sector)\\
$B = 12^{3/2}(2\eta(i))^{1/4} \approx 46.28$: \textbf{[NUM]} (numerical match only)\\
Identification $B=N_{\mathrm{eff}}^{3/2}\cdot Z_{1\mathrm{real}}^{1/4}$ from UBT action: \textbf{[OPEN/MC]}
\end{statusbox}

%──────────────────────────────────────────────────────────────────────────────
\section{Step 4: The Mellin Insertion — Exact Gap Statement}
\label{sec:mellin}
%──────────────────────────────────────────────────────────────────────────────

The single missing theorem required to close Gap G137-B is:

\begin{theorem}[Mellin insertion for $B$ — OPEN]
\label{thm:mellin}
Let $W_{\mathrm{eff}}$ be the one-loop effective action for the $N_{\mathrm{eff}} = 12$
UBT modes on $T^3_\psi\times S^1_\chi$ at the self-dual point $\tau = i$.
The coefficient $B$ in $V_{\mathrm{eff}}(n) = n^2 - Bn\ln n$ equals
\[
B = N_{\mathrm{eff}}^{3/2}\cdot Z_{1\mathrm{real}}(\tau=i)^{1/4}
\]
if and only if the effective action factorizes as
\begin{equation}
W_{\mathrm{eff}} = -N_{\mathrm{eff}}\,\mathrm{Tr}[\ln P_b]
\times N_{\mathrm{eff}}^{1/2},
\label{eq:weff_factor}
\end{equation}
where $P_b$ is the propagator for one charged $b$-field, and the
Mellin transform of $\mathrm{Tr}[\ln P_b]$ at the self-dual point gives
$Z_{1\mathrm{real}}^{1/4}$ upon summing over KK modes with divisor weights.
Status: \textbf{[OPEN]}.
\end{theorem}

\subsection*{What has been proved (up to this point)}

\begin{enumerate}
\item $Z'_{\mathbb{Z}^2}(0) = -4\ln|\Gamma(1/4)| + C_0$ via Chowla--Selberg [STD].
\item $n\ln n$ scaling from divisor sum asymptotics: [L1 asym.].
\item $\vartheta_3(0|i) = \sqrt{2}\,\eta(i)$ and $\Gamma(1/4) = 2\pi^{3/4}\eta(i)$: [STD].
\item $N_{\mathrm{eff}} = 12$ from SU(2)-twist: [L1].
\item Numerical value $B = 12^{3/2}(2\eta(i))^{1/4} \approx 46.28$: [NUM/OBS].
\end{enumerate}

\subsection*{What remains open}

The factorization \eqref{eq:weff_factor} from $S[\Theta]$ is not proved.
Specifically:
\begin{itemize}
\item The volumetric prefactor $N_{\mathrm{eff}}^{1/2}$ must be derived from
      the $T^3_\psi$ measure at self-duality; this is [OPEN].
\item The Mellin transform step $\mathrm{Tr}[\ln P_b] \mapsto Z_{1\mathrm{real}}^{1/4}$
      requires identifying $Z_{1\mathrm{real}} = 2\eta(i)$ from the UBT spectrum
      rather than setting it by hand; this is [OPEN/MC].
\item The connection between orbifold insertion
      $B_{\mathrm{orb}} = 12^{3/2}Z_{\mathrm{orb}}(i)^{1/4} \approx 44.40$
      and the target $B \approx 46.28$ requires an additional correction term;
      source of the $\approx 4\%$ discrepancy is [OPEN].
\end{itemize}

\subsection*{Python numerical verification}

The following script (see \texttt{tools/chowla\_selberg\_numerical.py}) verifies the modular values:

\begin{verbatim}
import math
# Known high-precision values (mpmath-validated)
# eta(i) = Gamma(1/4) / (2 * pi^(3/4))
Gamma14 = 3.6256099082219083  # Gamma(1/4)
pi = math.pi
eta_i = Gamma14 / (2 * pi**(3/4))  # = 0.768225422...
th3_i = math.sqrt(2) * eta_i       # = 1.086434811... (Ramanujan CM)

B_candidate = 12**1.5 * (2*eta_i)**0.25
# B_candidate = 46.2809...  [NUM]

# Consistency: 4*ln|Gamma(1/4)| from Chowla-Selberg
print(f"4*ln|Gamma(1/4)| = {4*math.log(Gamma14):.8f}")
# = 5.15209...
print(f"B candidate = {B_candidate:.8f}")
# = 46.28087...
\end{verbatim}

\begin{statusbox}
$Z'_{\mathbb{Z}^2}(0) = -4\ln|\Gamma(\tfrac14)| + C_0$: \textbf{[STD] — ověřit}\\
$B$ z průměru $V_{\mathrm{eff}}$ přes Dirichletovu asymptotiku: \textbf{[L1 asym.]}\\
$Z_{1\mathrm{real}} = 2\eta(i)$ z NS sektoru: \textbf{[OPEN/MC]} (not yet derived)\\
$B = N_{\mathrm{eff}}^{3/2}(2\eta(i))^{1/4} \approx 46.28$ analyticky: \textbf{[OPEN/MC]}\\
Volumetric prefactor $N_{\mathrm{eff}}^{1/2}$ from $T^3_\psi$ measure: \textbf{[OPEN]}\\
Gap G137-B: \textbf{[NARROWED — zbývá: exact Mellin/spectral insertion]}\\
Alpha: \textbf{NOT DERIVED}
\end{statusbox}

%──────────────────────────────────────────────────────────────────────────────
\section{Focused modular-bootstrap route (2-week timebox)}
\label{sec:modular_bootstrap}
%──────────────────────────────────────────────────────────────────────────────

To enforce a single strict route for Priority 3, we now run only the modular-bootstrap path:
\[
Z_{\mathrm{wind}}(\tau)=\vartheta_3(0|i\tau),
\]
with no input from $\alpha_{\mathrm{exp}}$ or fitted $B_{\mathrm{phenom}}$.

\subsection*{Step 1: modular differential structure}

We track the Ramanujan-type differential relation
\[
\vartheta_3'' = \left(\frac{\vartheta_3'}{\vartheta_3}\right)\vartheta_3'
+ \frac14\,E_2\,\vartheta_3,
\]
and evaluate at $\tau=i$, where $E_2(i)=3/\pi$ [STD].  This anchors the
self-dual point as the canonical bootstrap point for $Z_{\mathrm{wind}}$.

\subsection*{Step 2: bootstrap condition for $B$}

For
\[
V_{\mathrm{eff}}(n)=n^2-B\,n\ln n,\qquad
n_*=e^{B/2-1},
\]
the target fixed point $n_*=137$ implies
\[
B=2(\ln 137+1)\approx 46.28.
\]
Independently, the UBT modular candidate remains
\[
B_{\mathrm{mod}}(\tau)=12^{3/2}\,\bigl(2\eta(i\tau)\bigr)^{1/4}.
\]
The active test is whether modular constraints uniquely select a $\tau$ (in
particular $\tau=i$) that simultaneously gives the UBT $B$-value and the
stationary-point target under the same route assumptions.

\subsection*{Step 3: numerical scan protocol}

Numerical scan script: \texttt{tools/modular\_bootstrap\_scan.py}.  It computes
$B_{\mathrm{mod}}(\tau)$ and the fixed-point iteration estimate for $n_*(B)$ over
$\tau\in[0.5,2.0]$ to check uniqueness of the self-dual point.

\begin{statusbox}
Modular bootstrap route: \textbf{[aktivní — 2-week timebox]}\\
$B$ z modulárních podmínek: \textbf{[OPEN]}\\
Gap G137-B: \textbf{[NARROWED; go/no-go pending]}\\
Alpha: \textbf{NOT DERIVED}
\end{statusbox}

%──────────────────────────────────────────────────────────────────────────────
\section{Focused closure test: \(\tau=i\), \(W_{\mathrm{eff}}\), and verdict}
\label{sec:taui_focus}
%──────────────────────────────────────────────────────────────────────────────

\subsection*{Step 1 — Why \(\tau=i\)?}
\(\tau=i\) is the self-dual point of \(\mathbb{H}/\mathrm{SL}(2,\mathbb{Z})\),
fixed by \(S:\tau\mapsto-1/\tau\). Numerically, the active bootstrap scan reports
\(n_*(\tau=i)\approx136.99\), i.e. close to the target \(137\) [OBS].
This supports \(\tau=i\) as a preferred candidate point, but uniqueness at the
theorem level remains open in this note.

\subsection*{Step 2 — \(W_{\mathrm{eff}}\) factorization target}
Direct one-loop functional integration gives
\[
W_{\mathrm{eff}}=\det(P_b)^{-N_{\mathrm{eff}}/2},\qquad
\ln W_{\mathrm{eff}}=-\frac{N_{\mathrm{eff}}}{2}\,\mathrm{Tr}[\ln P_b].
\]
The missing closure step is the \(T^3_\psi\) volumetric lifting to the
effective \(N_{\mathrm{eff}}^{3/2}=N_{\mathrm{eff}}\cdot N_{\mathrm{eff}}^{1/2}\)
normalization in the \(B\)-coefficient map. This factorization remains the
core [OPEN] item.

\subsection*{Step 3 — Python numerical verification}
\begin{verbatim}
import mpmath, math
mpmath.mp.dps = 50

eta_i = float(mpmath.re(mpmath.eta(mpmath.mpc(0,1))))
B_target = 12**1.5 * (2*eta_i)**0.25
B_req = 2*137/(math.log(137)+1)

print(f"τ=i selektuje n*=137: [OBS]")
print(f"B_target = {B_target:.10f}")
print(f"B_req    = {B_req:.10f}")
print(f"Přesnost: {abs(B_target-B_req)/B_req*100:.6f}%")

# T-dualita argument
R_self_dual = 1.0  # v Planck jednotkách
print(f"R_self_dual = l_Pl = {R_self_dual}")
print(f"τ = i ↔ R·Λ = 1 ↔ self-dual pod T-dualitě")
\end{verbatim}

\begin{statusbox}
\(\tau=i\) jako T-duální self-dual bod: \textbf{[MC → L1 pokud T-dualita v UBT]}\\
\(n_*(\tau=i)=137\) analyticky: \textbf{[OBS → MC]}\\
\(W_{\mathrm{eff}}\) faktorizace \(N_{\mathrm{eff}}^{3/2}\): \textbf{[OPEN]}\\
Gap G137-B: \textbf{[NARROWED]}\\
Alpha: \textbf{NOT DERIVED}
\end{statusbox}

%──────────────────────────────────────────────────────────────────────────────
\section{Summary and Gap Closure Status}
\label{sec:summary}
%──────────────────────────────────────────────────────────────────────────────

\subsection*{Chain of reductions}

\[
\text{Gap G137-B}
\;\xrightarrow{\text{Step 1}}\;
Z'_{T^2}(0) \;\text{[STD]}
\;\xrightarrow{\text{Step 2}}\;
n\ln n \;\text{[L1 asym.]}
\;\xrightarrow{\text{Step 3}}\;
N_{\mathrm{eff}}^{3/2}\cdot Z_{1\mathrm{real}}^{1/4} \;\text{[NUM/MC]}
\;\xrightarrow{\text{Step 4 — OPEN}}\;
B \;\text{from } S[\Theta].
\]

\subsection*{What is needed for full closure}

\begin{center}
\begin{tabular}{lll}
\toprule
\textbf{Step} & \textbf{Claim} & \textbf{Status} \\
\midrule
1 & $Z'_{T^2}(0; d^2) = -2\ln R + Z'_{\mathbb{Z}^2}(0; d^2)$ & [STD] \\
1 & Chowla--Selberg: $Z'_{\mathbb{Z}^2}(0; d^2) = -4\ln|\Gamma(\tfrac14)| + 2\ln d + C_0$ & [STD] \\
2 & $n\ln n$ from Dirichlet divisor sum & [L1 asym.] \\
3 & $N_{\mathrm{eff}} = 12$ from SU(2)-twist & [L1] \\
3 & $B \approx 46.28$ numerically & [NUM] \\
4 & Factorization $W_{\mathrm{eff}} = N_{\mathrm{eff}}\cdot\mathrm{Tr}[\ln P_b]\times N_{\mathrm{eff}}^{1/2}$ & [OPEN] \\
4 & Mellin: $\mathrm{Tr}[\ln P_b] \mapsto Z_{1\mathrm{real}}^{1/4}$ with $Z_{1\mathrm{real}} = 2\eta(i)$ & [OPEN/MC] \\
\bottomrule
\end{tabular}
\end{center}

\subsection*{Next step}

Prove the factorization \eqref{eq:weff_factor} by computing $W_{\mathrm{eff}}$
directly from $S[\Theta]$ using the biquaternion field equations,
identifying the volume element on $T^3_\psi\times S^1_\chi$ at the
self-dual point, and performing the spectral sum via the Mellin transform
of the heat kernel $K_{P_b}(t) = \sum_\lambda e^{-\lambda t}$.

Related files:
\begin{itemize}
\item \texttt{research\_tracks/T3\_ALPHA/chowla\_selberg\_B\_derivation.tex}
  (Steps 1--7, predecessor of this note)
\item \texttt{canonical/alpha/alpha\_gap\_closure\_matrix.tex}
  (entries $\alpha25$--$\alpha26$)
\item \texttt{research\_tracks/quantum\_ubt/su2\_twist\_neff12.tex}
  ($N_{\mathrm{eff}} = 12$ [L1])
\item \texttt{tools/chowla\_selberg\_numerical.py}
  (Python verification of all numerical values)
\end{itemize}

\end{document}
